% \section{Contribution: Combining Knowledge Graphs and NLP to Analyze Instant Messaging Data in Criminal Investigations}
\section{WISE}
\todo{more than ee: exploit semi-structured chats for graph}
\todo{speak about rw on investigations in ee for legal above}

\todo{rw ima and investigations}

\todo{divide the contributions between extraction and access. merge all extraction and then all access}


It combines KGs to organize data extracted from IMA, deep learning-based speech-to-text technology to transcribe audio messages, NLP to extract entities and annotate chats, and semantic search to access the data.
Chat dumps include message content and metadata, such as participants, timestamps, and message types. The dumps are processed to populate a property graph, the chat KG, stored in Neo4j\footnote{https://neo4j.com/}. The graph represents relational data (chats, participants, exchanged messages) and message content. Multimedia enrichment generates textual content from audio messages and other attachments, enabling search functionality and entity extraction from this data. Currently, we focus on transcribing audio messages with the Whisper APIs~\cite{radford23a-whisper}, but the figure suggests potential extensions to image-to-text and video-to-text enrichment.
Entities in the chat metadata (senders, participants) are parsed to construct the graph. For message content, we employ a Named Entity Recognition and Linking (NEEL) pipeline adapted from previous work~\cite{ijckg2022,aixia_2023}.
It combines Named Entity Recognition (NER), Named Entity Linking (NEL), NIL prediction and entity clustering (for more details, see Section~\ref{sec:approach}).
With NEEL, we annotate the chat content, creating an annotated corpus, and we update the KG with the extracted entities.

In this paper, we consider two User Interfaces (UI) for accessing data and gaining insights:
the first is a native Neo4j UI that supports querying and visual exploration of the KG, which was a major factor in choosing Neo4j for storage. The second UI is DAVE (Document Annotation Validation and Exploration), a web application for visualizing, searching, and editing annotated documents, where users can use faceted search to filter results with entity-based facets. 

%\todo{spiegare come colleghiamo nlp e kg! nlp sfruttando info strutturate? triener?}

% \todo{feedback from stakeholders. mettere quando parliamo del progetto sora}

% Once messaging data is processed by the entity extraction pipeline...

% Our aim is to provide stakeholders the means to grasp knowledge out of these ``massive'' information sources.


% using entity extraction techniques, knowledge graphs, and recent speech-to-text models.

% With entity extraction techniques we refer to named entity recognition (NER), the task of \todo{}, named entity linking (NEL), \todo{}, NIL Prediction,, and NIL clustering. In fact these systems combined (their combination is often obtained with a pipeline) allow to extract the entities mentioned in textual documents, or in an entire chat.

% Then, knowledge graphs are suitable for representing connecting entities, such as chats, messages, people, locations, etc., so that these connections can be explored by end users. \todo{example queries}

% With respect to voice messages we use recent speech-to-text models to obtain the transcriptions to process with the entity extraction pipeline.

% Finally, while most work in legal entity extraction focused on NER, we believe that in a context such as investigations it is vital to go beyond mentions reaching a ``knowledge layer'' though the tasks of NEL, NIL prediction, and entity clustering and by representing the extracted knowledge in a knowledge graph.

%\todo{perchè facciamo nil prediction e clustering? cosa ne facciamo delle nuove entità}
%\todo{chiarire bene cosa facciamo già e cosa invece vogliamo fare}

%\todo{fig:1 caption}


% \todo{pensiamo a strutturare anche come sappiamo che andrebbe fatto:
% estrazione metadati in KG, ``multimedia enrichment'' (trascrizione audio, captioning) e ora lancio entity extraction (su un messaggio alla volta o su chat intera?) }



%\todo{un po' di soa. Recent advancement in ...}

\subsection{approach}

    % \subsection{Overview}

To provide more details about our approach, we describe the extraction of the knowledge graph (Section~\ref{sec:metadata_kg}), the methodology we adopt for multimedia (Section~\ref{sec:multimedia}), the entity extraction pipeline (Section~\ref{sec:entity_extraction}), and the exploratory interfaces (Section~\ref{sec:exploration}).
        % \begin{enumerate}
        %     \item what is a chat? what are the features of a message? 
        %     \item Processing of structured data: chats, messages, sender/receiver/people in groups.
        %     \item Entity extraction from message corpus and voice messages transcriptions.
        % \end{enumerate}
        % In a chat the information is divided between structured and unstructered data. Structured information include the participants of a chat and for each messages who the sender name (or the phone number), the receiver, and the time the message was sent. On the other hand, the message corpus and the transcription of a voice message come in the form of unstructured data.

        % For the former class of information, the structured data, we built a parser to extract the available structured information. Then we populate a knowledge graph using Neo4j creating the entities of chats, messages, people, and attachments.

        % For the latter, the unstructured, we employ our pipeline of entity extraction algorithms. At this point we
        
    \subsection{Optional Knowledge Graph} \label{sec:metadata_kg}
%Chats and chat messages, in addition to the unstructured corpus of the message, include structured data. 
Chat dumps are available in a copy of the mobile phone extracted with certified forensic applications. %and consist of textual files with some structure.
Information about each chat includes the list of participants with their phone numbers, the start time, and the time of the last activity. Each message is described by its timestamp, the name and number of the sender, and the attachments it contains.
We consider these metadata to be applicable to most instant-messaging applications. 
However, different forensic applications may export data in different formats. In our case, we received structured textual data for the first investigation and MS Excel (.xlsx) data for the second one.

% These metadata might depend on the specific IMA, e.g., WhatsApp, but do not depend on the forensic application used for the dump extraction. However different forensic application would likely export data in different formats.

% For the specific data format received in the experimented investigation, we build a parser based on parsing expression grammars (PEGs)~\cite{ford2004_peg}, using the parsimonious library\footnote{https://pypi.org/project/parsimonious/}, to extract all the aforementioned metadata and the content of the messages.

% While chats messages and metadata from the previous investigation were received in a structured textual form\footnote{For the previous investigation we used parsing expression grammars (PEGs) with parsimonious (https://pypi.org/project/parsimonious/).}, in the second investigation, the one we are considering in this work, the data received was in MS-XLSX (.xlsx), a machine readable format, easier to process.

Once extracted, metadata and messages contents are used to populate the Neo4j chat knowledge graph. A conceptual schema of the knowledge graph is shown in Figure~\ref{fig:model}; dashed arrows represent is-a relations, continuous arrows represent relations and properties. We also represent \textit{sameAs} relationships, which are extracted from the chat metadata (different numbers associated with the same contact).      
%shows the schema of the knowledge graph with the entity classes Chats, Messages, Person, and Attachment.%\todo{add attachment. add attributes.} 
The main reason to choose a property graph and Neo4j over semantic web standards is the opportunity to exploit its visual interface for reporting purposes.     

%Upon completing the metadata extraction, out of a total of 1501 chats, we successfully processed 1133 chats. This involved \todo{completa} individuals, \todo{completa} attachments, and \todo{completa} messages. (Queste informazioni sono già all'inizio della sessione 5)

%\todo{aggiungere visualizzazioni anonimizzate dei casi d'uso.}
    %\todo{numeri: 1501 chat WhatsApp di cui 1133 successi, quante persone nei metadati? quante entità?,... quanti allegati di vario tipo?}
    %\todo{entità}
    %\todo{attributi}
    % entita
    
%\todo{or metadata parsing}

\begin{figure}[ht]
\centering
\includegraphics[width=.8\textwidth]{images/wise/kg_model.drawio.pdf}
\caption{Knowledge Graph schema} \label{fig:model}
\end{figure}

\textbf{Knowledge graph exploration.}
Users can navigate the knowledge graph created by metadata using the query and visual exploration interface of Neo4j. This approach provides a comprehensive overview of the relationships and connections between various elements in the investigative data, but it also supports specific queries on the data. An example of this interface can be seen in Figure~\ref{fig:KGexploration}. The Figure shows all messages and chats between the Person Under Investigation (PUI) and another subject; sameAs connections reveal that the PUI and the other subjects have different numbers used in different chats.    

\begin{figure}
\includegraphics[width=\textwidth,trim=0 160 0 150, clip]{images/wise/KGexploration.png}
\caption{Analysis of the correspondences between PUI and a second individual% We provide the translation for relationships in Italian.
%The SameAs relationships help to identify the same person even when they use different phone numbers. Some nodes have been obscured due to the sensitivity of the information.
} \label{fig:KGexploration}
\end{figure}



\subsection{Exploration and Editing of Annotated Data} \label{sec:exploration}

%\todo{priorità max immagine con faceted search e chat}\todo{dobbiamo censurare tutto il testo presente nella chat con anche i nomi presenti nella faceted search}
%\todo{priorità immagine con trascrizione annotata}
%\todo{screenshot query demo modulo rag}



% The latter interface is in an early stage of development; it has been mostly introduced to collect feedback on the interest in this search interface from the investigators. In the rest of the paper, we focus on the first two interfaces in terms of technical details and evaluation, and only briefly discuss the third one.       

%\subsubsection{Knowledge graph exploration}


%DAVE provides a semantic retrieval system and a LLM-based question answering system; e.g., filtering the chats for a keyword such as ``withdraw'' and then using the entities as facets selecting a suspect. The former
%%
%%based on Elasticsearch\footnote{https://www.elastic.co/elasticsearch},
%%
%allows prosecutors to filter chats by keywords, metadata, and mentioned entities. The latter employs a generative LLM in a retrieval augmented generation (RAG) setting:
%%
%with a vector store (ChromaDB\footnote{https://www.trychroma.com/}) and an embedding model\footnote{https://huggingface.co/efederici/sentence-it5-base},
%%
%users can ask questions in natural language and receive answers with the reference to the documents containing the information they are asking about. 
%%
%% Once a chat is retrieved the user can open it and use the entity-driven exploration provided by DAVE. \todo{troppo specifico, non nell'intro}
%%
%For instance, a user question may be ``Is there evidence of any thefts? If so, what items were stolen?'' and the system would give an answer. Additionally, the system returns the chat messages that were used as the context to generate the answer, so that users can confirm the accuracy of the answer by referring to the original data. 
%%
%% Furthermore, in the case of voice messages, we maintain the association between transcribed words and the corresponding time in the audio, allowing users to listen to the specific segment without the need to go through the entire voice message or recording.
%%
%Indeed, LLM may produce hallucinated content~\cite{ji23hallucination}, that is nonsensical or unfaithful to the provided context.
%; in our example, an hallucination may happen when the model answers with ``Yes, there were several thefts and the stolen items were smartphones and wallets'' when no message in the context contained such information.

%\subsubsection{Faceted Search}
\textbf{Faceted Search and Document Exploration.}
Faceted search allows users to retrieve chats through a traditional keyword search. Users can filter results by selecting specific metadata (e.g., participants) or semantic annotations (e.g., other people, organizations, or places mentioned in the chat). Our current implementation is based on Elasticsearch\footnote{https://www.elastic.co/elasticsearch}. A peculiar feature is that semantic search is aware of links and clusters identified in the annotation process. 
A sanitized example of the Faceted Search interface is shown in Figure~\ref{fig:Faceted Search}; three chats containing the keyword ``incontro" (Italian for ``meeting") are found; these chats also contain specific persons and specific places shown in the panel on the left; by selecting a specific entity or metadata, the results are filtered. A user can click on a chat she/he wants to open using the \textbf{Document Explorer} page; from this interface, as shown in Figure~\ref{fig:document explorer}, the user sees the clusters of entities in the left panel and all their mentions; by clicking on a mention, she/he can move to the exact position of the mention. %\todo{screen}


% \begin{figure}
%     \centering
%     \includegraphics[width=0.8\textwidth]{Faceted Search.png}
%     \caption{Using the Faceted Search on DAVE to retrieve relevant information.
% %In this example, the messages are filtered to retrieve those citing a specific person and a specific place.
% Filters and chats have been obscured due to the sensitivity of the information.} 
%     \label{fig:Faceted Search}
% \end{figure}
% \vspace{-5pt}

% \begin{figure}
%     \centering
%     \includegraphics[width=0.8\textwidth]{document_explorer.png}
%     \caption{Using the Document Explorer to navigate the entity mentions in the chat.
% %In this example, the messages are filtered to retrieve those citing a specific person and a specific place.
% Sensitive information has been obscured.} 
%     \label{fig:document explorer}
% \end{figure}
% \vspace{-5pt}

\begin{figure}
    \begin{subfigure}[b]{0.49\textwidth}
        \centering
        \includegraphics[width=\textwidth, trim=0 0 270 0, clip]{images/wise/Faceted Search.png}
        \caption{Faceted Search}% Sensitive information has been obscured.} 
        \label{fig:Faceted Search}
    \end{subfigure}
    \hfill
    \begin{subfigure}[b]{0.49\textwidth}
        \centering
        \includegraphics[width=\textwidth, trim=0 0 100 0, clip]{images/wise/document_explorer.png}
        \caption{Document Explorer} 
        \label{fig:document explorer}
    \end{subfigure}
    \caption{DAVE interface} 
\end{figure}

%\centering
%\includegraphics[width=\textwidth]{Faceted Search.png}
%\caption{Utilization of Faceted Search on DAVE to retrieve relevant information.
%In this example, the messages are filtered to retrieve those citing a specific person and a specific place.
%Filters and chats have been obscured due to the sensitivity of the information.} \label{fig:Faceted Search}
%\end{figure}

% \subsubsection{Question Answering}
% The QA module works as a chatbot, allowing users to ask questions and receive answers. 
% %based on the knowledge extracted from the investigative data. 
% Users can explore the information through a conversational interface that utilizes a LLM within a retrieval augmented generation (RAG) framework. Messages (or chunks of messages) are embedded using an Italian T5 model\footnote{https://huggingface.co/efederici/sentence-it5-base}, then vectors are stored with ChromaDB\footnote{https://www.trychroma.com/}, a vector database, that is also responsible for the retrieval. The question is, indeed, vectorized by the embedder model, then related chunks are retrieved  by ChromaDB through approximate nearest neighbor. At this point the LLM (Llama 2 7B~\cite{touvron2023llama}) receives the question and the relevant messages in the input prompt and generates the answer. The answers is shown to the user maintaining a link to the original messages, enabling users to trace the source of the information with the ``Document Explorer''.
% %\todo{screen}

\subsubsection{Role of Different Interfaces and Annotation Editing}

While the search patterns supported by the two interfaces overlap, each interface has a specific role. 
Table~\ref{tab:interfaces role} summarizes the information each interface is able to provide the user, in which form, and the impact of the graph representation and of the NLP entity extraction to the results given to the user.
Additionally, DAVE was designed for a human-in-the-loop approach, considering the sensitive nature and ethical implications of the legal domain. Users can edit and correct annotations, ensuring the quality of the extracted knowledge.
%\todo{screen}   

%\todo{immagine workflow dell'allegato 4 dave ma adattata alle chat}

 % come la parte di processing impatta sulle interfacces


\begin{table}[h]
% \begin{tabularx}{\textwidth}{|X|X|X|X|X|}
\centering
\caption{Role of different interfaces}
% \begin{tabular}{|p{1.5cm}|p{1.5cm}|p{1.5cm}|p{3.5cm}|p{3.5cm}|}
\tymin=50pt
\tymax=200pt
\begin{tabulary}{\textwidth}{|L|L|J|J|J|}
% \begin{tabular}{|p{.12\linewidth}|p{.12\linewidth}|p{.12\linewidth}|p{.30\linewidth}|p{.30\linewidth}|} 
% \begin{tabularx}{\linewidth}{|X|X|X|X|X|} 
\hline
\textbf{Interface} &
  \textbf{Query} &
  \textbf{Result} &
  \textbf{Impact of Graph} &
  \textbf{Impact of NEEL} \\
  \hline
Graph &
  Cypher
  %, Simplified functions\tablefootnote{We developed specialized Python functions, e.g. matchMessages(keywords) or matchChats(kw), as a semplification of Cypher queries.} 
  &
  Tabular, Graph &
  Match all nodes attributes. Graph visualization &
  Queries match ``mentioned entity'' nodes \\
  \hline
% Document Interface &
DAVE &
  Faceted Search &
  Chats &
  Annotations linked to entities in the chat knowledge graph &
  % Given a set of chats that match a search,
  Entities in the matching chats are shown to the users and can be used to filter \\
  \hline
% Document Interface &
DAVE &
  Document Explorer &
  A specific position in a chat &
  Annotations linked to entities in the chat knowledge graph &
  Entities in a chat are shown to the users and can be used to filter\\
  \hline
% \end{tabularx}%
\end{tabulary}
\label{tab:interfaces role}
\end{table}

% \begin{table}[h]
% \small
% \begin{center}
% \begin{tabular}{|p | p | p| p| p|} 
%  \hline
%  \textbf{Interface} & \textbf{Query Type} & \textbf{Results} \\ [0.5ex] 
%  \hline\hline
%  Reconstruct relations between PUI and other persons  & Find all chats and messages between Suspect and Person 1   & Relational analysis  \\[1ex]
%  \hline
%  Search evidence of reported meeting between PUI and P1 & Find all messages that include one keywords in a list (e.g., “meeting”, “see you in”, etc.) & Text search  \\[1ex]
%  \hline
%  Search hints of money transfers between a list of People and Companies & Find all messages that cite an entity of type money  & Entity search \\[1ex]
%  \hline
%  Search for attachment types and distribution  & Find all attachment types and their respective frequency & Relational analysis \\[1ex]
%  \hline
%  People with whom PUI has interacted the most  & List of chats with highest number of interactions between PUI and different people & Relational analysis \\[1ex]
%  \hline
%  Chat naming a specific person  &Find all of the chat that cite an entity of type person & Entity search \\[1ex]
%  \hline
%  Persons with different phone numbers & Find the total number of phone numbers associated with the same person through SameAs relations  & Relational analysis \\[1ex]
%  \hline
% \end{tabular}
% \caption{Examples of queries on the graph}\label{Queries-on-graph}
% \end{center}
% \end{table}

\subsection{experiments}

\textbf{Impact of Entity-Based Modeling on Search.} 
In the first investigation, we used the tool to write a report. Therefore, we translated each investigative question into Neo4j queries and reported the findings, providing: the queries, arguments and explanations for their choice, graph-based visualizations provided by the UI (e.g. Figure~\ref{fig:KGexploration}), and copies of the relevant messages. We can classify the query patterns into three main groups based on their purpose: relational analysis (e.g., between PUI and other people), text search, and entity search.  
%Examples of investigative questions, queries, and their category, are reported in Table~\ref{Queries-on-graph}. 
%Please note that these are some examples of questions and queries.
% Multiple queries were constructed for each investigative question.

We found that keyword and entity searches were very frequent. After validating the utility of combining graph-based constraints and textual search, which allows visualizing nodes (messages) containing the query pattern and exploring the connected nodes, and including graph-based visualizations in reports, we further explored the impact of our entity-centric approach in the second investigation, where messages and transcripts increased significantly. Law enforcers provided a list of 30 keywords of interest (22 generic terms and 8 entity names) for evaluation.
The total results from searching each keyword amounted to 818 chats, 1780 messages, and 230 audio transcripts (results for different keywords may overlap). The generic terms (translated in English) returning the most results (chats/messages) are ``agreement'' (126/274), ``runoff'' (113/133), and ``tender'' (92/296). These numbers indicate a high results space for keyword searches, highlighting the need for advanced mechanisms to filter and explore data with specific entity constraints.
For each query, we counted the distinct entities present in the retrieved chats, distinguishing between entities from IMA metadata and those recognized by the NEEL pipeline from the message corpus, which are more prone to errors. The median number of distinct entities retrieved from chats metadata is 181, and 568 for entities recognized by the NEEL pipeline. These entities can filter the often large number of results (e.g., consider only messages sent by X or in chats with Y) or provide insights into entities relevant to the query context. The data volume and associations between entities and results support the introduction of faceted search and document explorer interfaces, which display the entities found in the results.

%the median number of distinct entities associated in the graph with the chats returned for a search result is 181.  returned search result in the graph (without counting entity mentions) is 181, while  occurring in every chat returned in the results 
The terms referring to entities (persons and organizations) return fewer results (max 56 chats, min 0). Two terms are found in audio transcripts, providing evidence for the impact of the multimedia enrichment. We also match these terms to names of distinct entities extracted from the messages' content with the NEEL pipeline or included in the graph (w/o considering the NEEL pipeline). While three entities are found in both sources, four are found in only one source, and one in none. We believe that this observation suggests that combining keyword search with filters based on entities identified with two criteria (rule-based graph construction and NEEL-based extraction) is beneficial.       

%returning more results  for generic terms. These numbers suggest that advanced mechanisms to filter out results and explore the data are needed to get      chats are  Searching for these keywords in the graph return match jointly match  
%To support our main claim that entity-based filters  graph-based constraints are provide evidence for the need of      

%In particular, the keyword search was the most frequently used query example for finding messages or audio transcripts of interest. In particular, the keyword search directly on the graph, through the use of simple queries, makes it possible to visualise the nodes (messages) that contain the indicated word and to navigate the graph, allowing the nodes connected to the one of interest, such as the chat, to be displayed. 
%In addition, it was possible to identify a very limited number of messages containing words not in common use, such as names of particular organisations or nicknames of people, which would have taken time or would have been difficult to identify as they are details that may escape a superficial first reading. 

% \textbf{Statics Investigative questions and graph-based queries.} ??
% In the chats containing messages, the NEEL pipeline found 7486 mentions of entities: 4592 of type "person", 1166 of type "location", 606 of type "date", 229 of type "money" and 859 of type "organization". We should note that the addition of voice-messages transcriptions to the documents helped to recognize 138 more entities than when the pipeline was applied to the text messages alone. This suggests that voice messages may have a higher information content compared to plain-text messages, since the average number of entities mentioned in them (0.24) is higher than the average number of mentions in a text-message (0.16).\\

\textbf{Investigators' Feedback}. The feedback on our work was collected based on the report and demonstrations of our PoC (live demo) across the collaboration stages.
Investigators found the graph-based visualizations helpful and understandable. They are familiar with these structures because of the relevance of relational knowledge in their investigations. The combination of text search with relational queries was appreciated. In relation to multimedia data enrichment, they believe that speech-to-text could be really transformative in terms of speed and effort balance.
They have been quite surprised by the quality of the transcriptions and not concerned by occasional mistakes, as they always verify the evidence against the original audio. 
The NLP part was also appreciated in terms of potential and usefulness. The opportunity to quickly browse through the chat exploiting DAVE interface was deemed promising because their previous tools only supported syntactic match patterns.           

\textbf{Limitations and Challenges.}
%
%The NEEL stack should be better customized for the domain, e.g., via in-domain fine-tuning. %Clustering is attracting the attention of the research community and we believe that investigations and IMA data provide a representative and challenging domain to develop better algorithms for this task. 
%\todo{forse potremmo togliere questa discussione sul clustering visto che non è tanto esplorato nel paper}
%In the customization of entity clustering in the domain of IMA, we also need to improve the interconnection between the knowledge graph and the NEEL pipeline.
%For this reason, we can't yet provide any additional statistics on the number of different ``entity" nodes.\todo{Per Prof: controllare se OK}
% The integration of generative LLMs is still limited.
% The QA module requires significant improvement and
%We did not include image-to-text enrichment yet, a highly desired features.
After our demonstrations, investigators were eager to use our solution, but querying Neo4j through its interface requires expertise, and expressed their interest in conversational search interfaces to answer questions. We therefore developed two additional prototype interfaces 1) a set of intuitive and composable Python functions that handle a wide range of common query patterns,
%to simplify analyses based on the graph data, 
and 2) a conversational search interface based on the popular paradigm of retrieval augmented generation. We did not describe these features in this paper as their development stage is too early. However, the prosecutor and the law enforcers actually accessed the simplified query interface and were able to submit queries.    \todo{mention the simplified language? maybe not}
%is too early to be described in this paper have been shown and tested by our users (for example  to collect some initial feedback but their development stage is too early to be described in this paper. 